% !TEX root = ../thesis.tex
\chapter{Theoretical background} \label{sec:theorie}
\section{Magnetism}\label{sec:theory_magnetism}
\subsection{Atomic magnetism}\label{sec:theory_atomic_magnetism}
\subsection{Magnetic interactions}\label{sec:theory_magnetic_interactions}
\subsection{Thin film magnetism}\label{sec:theory_thin_film_magnetism}
\subsection{Magnetic particles}

\section{Magnetic anisotropy}\label{sec:theory_magnetic_anisotropy}
\subsection{Magnetocrystalline anisotropy}
\subsection{Shape anisotropy}
\subsection{Surface- and interface anisotropies}


\section{Exchange bias}\label{sec:theory_exchange_bias}

\section{DLVO-Forces}\label{sec:theory_dlvo}
\subsection{Electrostatic intercations}
\subsection{Van der Waals}

\section{Traveling wave magnetophoresis}
In this section, the theoretical background of traveling wave magnetophoresis and its implications are explained.
We will start with ...

\subsection{Particle transport mechanism}
\subsection{Close to surface transport}

$\#TODO$ where to put Reynolds number?
\begin{equation}\label{eq:Reynolds_number}
    \mathrm{Re} = \frac{\rho u L}{\mu},
\end{equation}
where at room temperature, $rho=\SI{1000}{\kg\per\cubic\meter}$ is the density of water, $u$ is the flow speed which, for our system is the velocity of the \acp{MP} and is generously given by \SI{100}{\micro\m\per\s}, $L$ is the charachtersitic length which is $<\mathrm{cm}$ range, and $\mu$ is the dynamic viscosity, given at \SI{1e-3}{\Pa}. 
It becomes apparent that at small volumes used, in the present case in the \SI{200}{\micro\L} range, $\mathrm{Re}$ becomes $<1$\cite{Shanko2019}.

Since the volumes used in this \ac{LOC} device are in the sub-mL range,  
Using relatively weak magnetic fields in the range of \SIrange{1}{4}{\kilo\A\per\m}, the superparamagnetic bead's magentic moment $\mathbf{m_\mathrm{MP}}$ depends on the external magnetic field linearly as 
\begin{equation}
    \mathbf{m_\mathrm{MP}} = \bar{\chi}_\mathrm{MP} V \mathbf{H}, 
\end{equation}
where $V$ is the bead's volume, and $\bar{\chi}=\nicefrac{3\chi_\mathrm{MP}}{\chi_\mathrm{MP} + 3}$ is the shape corrected bead susceptibility\cite{Yellen2007}. 
It becomes clear that the \ac{MP} is treated as a point dipole, which is a reasonable approximation when the magnetic field does not vary drastically across the \ac{MP}'s volume, but will be discussed later \cite{Yellen2007}.


When the \ac{MP} sees the magnetic field, the potential energy $U_\mathrm{P}$ of the \ac{MP} at lateral position $x$ and above the substrate $z$ is given by\cite{Holzinger2015a} 
\begin{equation}\label{eq:U_P_MP}
    U_\mathrm{P}(x,z) = -\mu_0 \mathbf{m_\mathrm{MP}}(x,z) \mathbf{H}(x,z), 
\end{equation}
where $\mu_0$ is the vacuum permeability.
It becomes apparent that $U_\mathrm{P}(x,z)$ is smallest where $\mathbf{H}(x,z)$, which is where the magnetic stray fields emerge from the substrate, that is, at the \acp{DW}.

External fields...

The magnetic force exerted on the superparamagnetic bead in the direction of motion $x$ is given by 
\begin{equation}
    F_\mathrm{m,x}(x,z,t)=V(\chi_\mathrm{MP}-\chi_\mathrm{fl}) \mu_0  ((\mathbf{H}_\mathrm{eff} \mathbf{\nabla}) \mathbf{H}_\mathrm{eff})_\mathrm{x},
\end{equation}
where $\chi_\mathrm{fl}$ is the susceptibility of the fluid, and $\mathbf{H}_\mathrm{eff}$ is the effective field, which is the sum of the static stray field emerging from the \ac{hh}-\ac{tt}-\ac{DW}, and superposed external fields created by the Helholz coils, at the position of the particle\cite{Shubbak2026}.
Since the \ac{MP} is residing in the potential minimum in accordance with \ref{eq:U_P_MP}, shifting this energy well will force the \ac{MP} to follow its motion. 
Fittingly, the transport concept is coined Traveling Wave Magnetophoresis CITE WHO?


Using an equation of motion, the transport from one \ac{DW} to the adjacent by periodically transforming $U_\mathrm{P}(x,z)$ can be described by\cite{Yellen2007, Shubbak2026}
\begin{equation}
    M \frac{d^2x}{dt^2}-f_\mathrm{D}\frac{dx}{dt} = F_\mathrm{m,x} \sin(\phi),
\end{equation}
where $M$ is the mass of an \ac{MP}, $f_\mathrm{D}$ is the viscous drag coefficient, and the phase $\phi = kx-\omega t$ denotes the relative position of the \ac{MP} with respect to the position of the nearest energy minimum. 